\chapter{Jaw Dislocation}

\section*{Definition and Overview}
Jaw dislocation occurs when the lower jaw (mandible) slips out of its joint with the skull, the temporomandibular joint (TMJ). It usually happens when the mouth is opened very wide, such as during yawning, dental treatment, vomiting, or a blow to the face, and it is more likely in people with lax joints, a history of dislocation, or TMJ problems. A dislocated jaw cannot close the mouth, is very painful, and distorts the face, and the person cannot speak or swallow normally. It is treated by a doctor or dentist repositioning the jaw, often with sedation or anaesthesia, and most people recover fully, though recurrence is possible.

\section*{Epidemiology}
Jaw dislocation is less common than dislocations of the shoulder or fingers but is well recognised in emergency departments. It can occur at any age, but is most common in young adults and in people with hypermobile joints, and it is slightly more common in women. Risk factors include a history of dislocation, TMJ disorders, dental procedures, seizures, and connective tissue laxity. Recurrent dislocation affects a small proportion of people and may need further treatment. Bilateral dislocation, affecting both joints, is more common than dislocation on one side.

\section*{Aetiology and Pathophysiology}
The temporomandibular joint is a hinge-and-slide joint in front of each ear, formed by the condyle of the jaw and a socket in the skull, separated by a disc. The jaw normally dislocates when the condyle slides forward beyond its socket and the muscles clamp it there, so it cannot return. This happens when the mouth opens beyond its normal range, as in yawning, dental work, laughing, or trauma, particularly when the muscles are relaxed and the joint is lax. Once dislocated, the jaw is locked with the mouth open, the chewing muscles go into spasm, and the position causes pain and difficulty swallowing or speaking. Repeated episodes can stretch the joint capsule and cause habitual dislocation.

\section*{Clinical Features}
\begin{itemize}
    \item Inability to close the mouth fully
    \item Pain in front of the ears and in the jaw
    \item A visible or palpable depression in front of the ear where the joint is empty
    \item The jaw may be shifted to one side in unilateral dislocation
    \item Difficulty speaking, swallowing, and managing saliva
    \item Facial muscle spasm and drooling
    \item A history of yawning, dental treatment, or trauma
    \item In recurrent cases, episodes that reduce spontaneously or with simple techniques
\end{itemize}

\section*{Differential Diagnosis}
Other jaw and facial problems should be distinguished from dislocation.
\begin{itemize}
    \item \textbf{Fracture of the jaw or cheekbone:} after trauma, with swelling and deformity
    \item \textbf{TMJ dysfunction:} clicking, pain, and limited opening without full dislocation
    \item \textbf{Locked jaw from muscle spasm:} trismus without joint displacement
    \item \textbf{Tetanus:} muscle rigidity with difficulty opening the mouth
    \item \textbf{Temporomandibular joint arthritis or infection:} pain and stiffness
    \item \textbf{Angioedema and other causes of facial swelling:} with difficulty swallowing
\end{itemize}

\section*{When to Seek Medical Care}
A dislocated jaw is a medical emergency because the person cannot close the mouth, swallow, or manage saliva. Go to an emergency department, where the jaw can be repositioned with sedation or anaesthesia. Anyone with a jaw injury that has caused dislocation, deformity, or an inability to bite normally should be assessed, and X-rays are often needed to exclude fracture.

\textbf{Call 000 (or local emergency services) for:}
\begin{itemize}
    \item A jaw that will not close after dislocation or injury
    \item Difficulty breathing or swallowing
    \item Heavy bleeding from the mouth after jaw injury
    \item Numbness, severe swelling, or deformity of the face after trauma
    \item Loss of consciousness or confusion after a facial injury
\end{itemize}

\section*{Diagnosis and Investigations}
\begin{itemize}
    \item \textbf{History and examination:} the mechanism, symptoms, and the position of the jaw
    \item \textbf{X-rays:} to confirm dislocation and exclude fracture
    \item \textbf{Panoramic dental X-ray:} good views of the jaw joints
    \item \textbf{CT scan:} for complex injuries, fractures, or recurrent problems
    \item \textbf{Neurological assessment:} for injury to the facial nerve after trauma
    \item \textbf{Assessment of recurrent dislocation:} dental and specialist review
\end{itemize}

\section*{Management}
\begin{itemize}
    \item \textbf{Repositioning:} the jaw is guided back into place by a clinician, usually with sedation or local anaesthesia to relax the muscles
    \item \textbf{Muscle relaxation:} medicines to allow the jaw to slip back
    \item \textbf{Jaw support:} a bandage or soft diet for a few days to allow recovery
    \item \textbf{Pain relief:} simple analgesics and anti-inflammatories
    \item \textbf{Avoiding wide opening:} support the jaw when yawning and during dental work
    \item \textbf{Physiotherapy:} for muscle and joint rehabilitation
    \item \textbf{Recurrent dislocation:} injections, physiotherapy, or rarely surgery to stabilise the joint
    \item \textbf{Treatment of underlying TMJ problems:} dental and specialist management
\end{itemize}

\section*{Complications}
\begin{itemize}
    \item Recurrent dislocation, which can become habitual
    \item Damage to the joint, disc, or joint capsule
    \item Fracture of the jaw condyle or skull base in trauma
    \item Injury to the facial nerve or nearby structures
    \item Difficulty eating, speaking, and maintaining oral hygiene
    \item Aspiration from inability to manage saliva
    \item Chronic TMJ pain and dysfunction
\end{itemize}

\section*{Prognosis and Outcomes}
A single jaw dislocation usually reduces easily and heals without long-term problems. Recurrence is the main concern, affecting up to a third of people, and can be managed with physiotherapy, behavioural techniques, and in persistent cases, injection therapy or surgery. Most people return to normal jaw function within weeks. Preventing over-opening and treating TMJ problems reduce the risk of further episodes.

\section*{Prevention and Self-Care}
\begin{itemize}
    \item Support the chin when yawning to limit mouth opening
    \item Tell your dentist about a history of dislocation before dental treatment
    \item Avoid very wide biting into large foods
    \item Practise TMJ exercises if advised
    \item Use relaxation techniques to reduce jaw clenching
    \item Seek early treatment for recurrent dislocation
    \item Follow soft-food and activity advice after reduction
    \item Wear mouthguards for contact sport if advised
\end{itemize}

\section*{Sources and Further Reading}
The following reputable sources provide further information for healthcare students and patients seeking reliable, current advice about jaw dislocation.
\begin{itemize}
    \item Healthdirect Australia. \textit{Jaw dislocation.} https://www.healthdirect.gov.au/
    \item Australian Dental Association. \textit{TMJ and jaw joint information.} https://www.ada.org.au/
    \item Better Health Channel. \textit{Temporomandibular joint disorders.} https://www.betterhealth.vic.gov.au/
    \item NHS. \textit{Dislocated jaw.} https://www.nhs.uk/
    \item Mayo Clinic. \textit{Temporomandibular joint disorders.} https://www.mayoclinic.org/
    \item MSD Manual. \textit{Dislocation of the jaw.} https://www.msdmanuals.com/
\end{itemize}
